\subsection{有限域代数闭包的伽罗瓦群}

设 \(S \neq \{1\}\) 是一组整数 \(\geqslant 1\) 在乘法下稳定；我们将通过关系“ \(m\) 整除 \(n\) ”来排序它。当 \(m\) 整除 \(n\) 时，有 \(mZ \supset nZ\) ，因此存在典范同态 \(\pi_{m,n}\) ，它把环 \(Z/nZ\) 映到环 \(Z/mZ\) 上。我们用 \(A(S)\) 表示环的逆向系统 \((Z/mZ, \pi_{m,n})\) （以 \(S\) 为指标集）的逆向极限。每个有限集 \(Z/mZ\) 都赋予离散拓扑，而 \(A(S)\) 赋予由乘积 \(\prod_{n \in S} (Z/nZ)\) 的拓扑诱导的拓扑。

于是 A(S) 是一个紧拓扑环（Gen.～Top. I，第64页，命题8）。我们立即看出，唯一的环同态 \(\varphi\) （由 \(\mathbf{Z}\) 映入 A(S) ）是单射且具有稠密的像；我们将把 \(\mathbf{Z}\) 与它在 \(\varphi\) 下的像在 A(S) 中等同。对于由 A(S) 的拓扑在 \(\mathbf{Z}\) 上诱导的拓扑，集合 \(m\mathbf{Z}\) （ \(m \in S\) ）构成 0 的邻域基。

当 \(S = N^{*}\) 时， \(A(S)\) 记作 \(\hat{Z}\) 。当 S 由一个素数 l 的幂组成时， \(A(S)\) 记作 \(Z_{l}\) ，称为“l 进整数环”。因此我们有

\[
\hat {\mathbf {Z}} = \varprojlim_ {m \geqslant 1} \mathbf {Z} / m \mathbf {Z}, \quad \mathbf {Z} _ {l} = \varprojlim_ {n \geqslant 0} \mathbf {Z} / l ^ {n} \mathbf {Z}.
\]

当 S 和 T 是两个在乘法下稳定的整数集合且满足 \(S \supset T\) 时，我们有一个自然投影 \(\mathbf{A}(\mathbf{S}) \to \mathbf{A}(\mathbf{T})\) ，它是拓扑环的连续同态。特别地，对每个素数 l，我们有一个连续同态 \(\hat{Z} \to Z_{l}\) 。由此得到连续同态

\[
\hat {\mathbf {Z}} \rightarrow \prod_ {l} \mathbf {Z} _ {l}
\]

（乘积取遍所有素数）；它是拓扑环的同构，这由对 I，第112页，命题11 取逆向极限得出。

设 K 是一个有 q 个元素的有限域， \(\bar{K}\) 是 K 的一个代数闭包。对每个整数 \(m \geqslant 1\) ，我们用 \(K_{m}\) 表示 \(\bar{K}\) 中在 K 上次数为 m 的唯一子域（命题3）。我们有 \(\bar{K} = \bigcup_{m \geqslant 1} K_{m}\) 。此外，我们用 \(\sigma_{q}\) 表示完备域

\(x \mapsto x^{q}\) 所定义的 \(\bar{K}\) 的自同构；它称为 \(\bar{K}\) 的（相对于 K 的）弗罗贝尼乌斯自同构。

命题5. --- 存在唯一的拓扑群同构 \(\pi_{\mathbf{K}}: \hat{\mathbf{Z}} \to \operatorname{Gal}(\bar{\mathbf{K}} / \mathbf{K})\) 使得 \(\pi_{\mathbf{K}}(1) = \sigma_q\) 。

设 \(\Gamma\) 为 \(\operatorname{Gal}(\bar{\mathbf{K}} / \mathbf{K})\) 中由 \(\sigma_q\) 生成的子群。对每个整数 \(m > 0\) ，我们有 \(\sigma_q^m (x) = x^{q^m}\) 对所有 \(x\in \bar{\mathbf{K}}\) 成立，因此 \(\sigma_q^m\) 的不动点集等于 \(\mathbf{K}_m\) 。由于 \(\mathbf{K}_m\neq \bar{\mathbf{K}}\) ，故 \(\sigma_q^m\neq 1\) 。因此存在同构 \(\pi_0\) ，它把 \(\mathbf{Z}\) 映到 \(\Gamma\) 上，并将 1 映到 \(\sigma_q\) 。

\(\Gamma\) 的不变量域由 \(x \in \bar{K}\) 中使 \(x^{q} = x\) 成立的那些 x 组成，因此等于 K。于是（V，第67页，引理2）群 \(\Gamma\) 在 \(\operatorname{Gal}(\bar{\mathbf{K}}/\mathbf{K})\) 中稠密。由于 \(\bar{K}\) 在 K 上的每个有限次子扩张都是诸域 \(K_{m}\) 之一，子群 \(\operatorname{Gal}(\bar{\mathbf{K}}/\mathbf{K}_{m})\) 构成 \(\operatorname{Gal}(\bar{\mathbf{K}}/\mathbf{K})\) 中 1 的基本邻域系。显然 \(\Gamma \cap \operatorname{Gal}(\bar{\mathbf{K}}/\mathbf{K}_{m})\) 是由 \(\sigma_{q}^{m}\) 生成的循环群，因此等于 \(\pi_{0}(m\mathbf{Z})\) 。

根据上述关于 \(\hat{Z}\) 的拓扑的论述，同构 \(\pi_{0}:Z\to\Gamma\) 以唯一的方式延拓为拓扑群的同构 \(\pi_{K}:\hat{Z}\to\operatorname{Gal}(\bar{\mathbf{K}}/\mathbf{K})\) 。

设 \(m \geqslant 1\) 为一个整数；显然， \(\bar{\mathbf{K}}\) 相对于 \(\mathbf{K}_m\) 的弗罗贝尼乌斯自同构就是 \(\sigma_q^m\) 。由此得到关系式

\[
\pi_ {\mathrm{K} _ {m}} (a) ^ {*} = \pi_ {\mathrm{K}} (m a) \quad \text { for } \quad a \in \hat {\mathbf {Z}}.\tag{5}
\]

